\chapter{Theoretical Background}
\label{chap:ch3}

In this chapter, we present the theoretical foundations underlying the proposed approach for 3D fractured object reassembly. We begin by introducing point cloud representations and their unique structural properties, which motivate the need for specialized processing techniques. We then review key deep learning architectures for point sets, including PointNet and PointNet++, highlighting their strengths and limitations in capturing geometric features. Building on this, we introduce transformer-based models and attention mechanisms, which enable richer modeling of relationships between points and across fragments. Finally, we describe the RANSAC algorithm and its integration with Singular Value Decomposition (SVD) for robust estimation of rigid transformations, providing the mathematical tools necessary for accurate geometric alignment.

\section{Point Cloud Representation}

A common representation for 3D geometry in learning-based pipelines is the point cloud, usually chosen due to their simplicity and flexibility. A point cloud is an unstructured set of 3D points, where each point \(p_i\) may include additional attributes such as color or surface normals. Unlike volumetric voxel grids, which suffer from cubic memory costs and resolution limits, or 3D meshes, which introduce complex and irregular topological connectivity, point clouds provide a highly efficient raw geometric representation. However, effectively processing these raw spatial data requires neural network architectures to explicitly account for three intrinsic structural properties:
\begin{itemize}
    \item \textbf{Permutation invariance;} Because point clouds lack a rigid spatial grid, the physical shape they represent remains identical regardless of how the points are ordered in memory. Consequently, any model processing them must be permutation invariant, ensuring it produces the exact same output regardless of the input data sequence.
    \item \textbf{Local geometric structure;} A single isolated point provides no meaningful contextual information. Identifiable geometric features, such as fracture edges, corners, or smooth surfaces, emerge only from the spatial proximity and interaction between neighboring points. Models must therefore be capable of grouping points to capture these localized structures.
    \item \textbf{Rigid transformation invariance;} The intrinsic geometry of an object does not change based on its location or orientation. Therefore, the learned representations must remain consistent under global rigid transformations, such as 3D translations and rotations in space.
\end{itemize}

\begin{figure}[h]
    \centering

    \begin{subfigure}[b]{0.32\textwidth}
        \centering
        \includegraphics[width=\textwidth]{figures/mesh.png}
        \caption{Mesh}
    \end{subfigure}
    \hfill
    \begin{subfigure}[b]{0.32\textwidth}
        \centering
        \includegraphics[width=\textwidth]{figures/pointcloud.png}
        \caption{Point Cloud}
    \end{subfigure}
    \hfill
    \begin{subfigure}[b]{0.32\textwidth}
        \centering
        \includegraphics[width=\textwidth]{figures/voxel.png}
        \caption{Voxel Grid}
    \end{subfigure}

    \caption{Common 3D shape representations. (a)~A triangular mesh encodes surface topology through vertices and faces. (b)~A point cloud represents geometry as an unordered set of 3D coordinates. (c)~A voxel grid discretizes 3D space into a regular volumetric lattice.}
    \label{fig:3d_representations}
\end{figure}

Because of these unique properties, standard Convolutional Neural Networks are fundamentally ill-equipped to process point clouds directly (Figure~\ref{fig:3d_representations}). Classical convolutions are strictly based on the fixed spatial order and uniform density of regular grids. Since point clouds are unordered sets scattered continuously in 3D space, traditional convolutional operations cannot be applied, necessitating the use of specialized, order-agnostic architectures.

\section{PointNet and PointNet++}
\label{sec:pointnet}

To address the inherent challenges of processing unstructured 3D data, early research introduced architectures designed specifically for point sets, like PointNet \cite{qi2017pointnet}. Instead of aggregating information from local neighbors using a spatial kernel and extracting a feature from a local patch of pixels, PointNet learns features for each point independently. It replaces the convolutional kernel with a shared Multi-Layer Perceptron (MLP), denoted \(h:\mathbf{R}^3 \xrightarrow{} \mathbf{R}^K\), which acts as a pointwise feature extractor. For every individual point \(x_i\), the network computes a high-dimensional feature vector \(h(x_i)\). 

But because these features are extracted in complete isolation, the network requires a mechanism to capture the global structure of the object while remaining strictly invariant to the permutation of the input points. PointNet achieves this through a symmetric aggregation function, specifically global max pooling:
\begin{equation}
\label{eq:pointnet_pool}
g=\max_{i=1 \ldots n}h(x_i)
\end{equation}
This works because \(\max(a, b)=\max(b, a)\). So, while CNNs build global understanding progressively by combining localized structures, PointNet defines global understanding as a collected set of the most dominant pointwise responses (Figure~\ref{fig:pointnet}). 

\begin{figure}[h]
    \centering
    \includegraphics[width=1\linewidth]{figures/pointnet.png}
    \caption{The PointNet architecture \cite{qi2017pointnet}. Each point is independently transformed by a shared MLP, and a symmetric max-pooling operation aggregates the per-point features into a single global descriptor.}
    \label{fig:pointnet}
\end{figure}

While the foundational PointNet formulation successfully guarantees permutation invariance through global max-pooling, it introduces a critical limitation: by design, it does not capture the local structures such as spatial distances and geometric interactions. To overcome this lack of local shape awareness, PointNet++ \cite{qi2017pointnetplusplus} introduces a hierarchical architecture that applies PointNet recursively on a nested partitioning of the input point set. By exploiting the underlying metric space, this hierarchy allows the network to capture fine-grained local features at continuously increasing contextual scales, analogous to the progressively larger receptive fields found in CNNs.

This hierarchical feature learning is achieved through an encoder made of repeated Set Abstraction levels. Each Set Abstraction level progressively downsamples the point cloud by using Farthest Point Sampling \cite{qi2017pointnetplusplus} to select a sparse set of centroids, groups neighboring points within a fixed metric radius, and applies a miniature PointNet to encode these localized regions into highly descriptive feature vectors. Because real-world point clouds often suffer from non-uniform sampling densities, PointNet++ also employs density-adaptive strategies, such as Multi-Scale Grouping, which intelligently combine features extracted from neighborhoods of varying sizes to ensure robustness (Figure~\ref{fig:pointnet++}).

\begin{figure}[h]
    \centering
    \includegraphics[width=1\linewidth]{figures/pointnet++.png}
    \caption{The PointNet++ architecture \cite{qi2017pointnetplusplus}. Hierarchical Set Abstraction levels progressively downsample the point cloud, applying local PointNet modules at each scale to capture fine-grained geometric features with increasing contextual scope.}
    \label{fig:pointnet++}
\end{figure}

\section{Transformers and Attention Mechanisms}

Although PointNet++ successfully captures local geometric structures using hierarchical neighborhoods, its reliance on max-pooling introduces some limitations. Max-pooling identifies the most prominent feature within a local region but discards the complex geometric relationships between the remaining points. For tasks that require identifying compatible fracture boundaries, the network must understand not just the presence of a feature, but exactly how different points relate to one another. To overcome this, modern architectures replace max-pooling with the Attention mechanism \cite{vaswani2017attention}, the foundational building block of Transformer networks.

Originally designed for natural language processing to model long-range dependencies between words in a sentence, transformers are remarkably well-suited for 3D point cloud processing. Fundamentally, the self-attention operator at the core of a transformer is a set operator. Because it computes relationships by comparing elements against one another rather than relying on a fixed grid or sequential order, it is invariant to both the permutation and cardinality of the input data. This makes the transformer a perfect mathematical match for processing unordered point clouds.

At its core, the attention mechanism can be described as mapping a Query and a set of Key-Value pairs to an output representation. The features of each point are linearly projected into three distinct vectors:

\begin{itemize}
    \item Query (Q): Represents what a specific point is looking for
    \item Key (K): Represents what a point has (its own structural identity)
    \item Value (V): The actual information of the underlying feature of that point
\end{itemize}

To compute the relationships between points, the network calculates the dot product between the Query of one point and the Keys of all other points. A high dot-product score indicates strong geometric compatibility or relevance. These scores are passed through a softmax function to create normalized attention weights, which dictate exactly how much of each point's Value should be mixed into the final aggregated feature. Instead of just taking the maximum value like PointNet, attention creates a dynamically weighted sum of the entire neighborhood based on structural context.

In the context of 3D fractured object reassembly, architectures like Jigsaw \cite{lu2023jigsaw} utilize two distinct variations of this mechanism to establish geometric matches:

\paragraph{Self-Attention.} In this layer, the Queries, Keys, and Values all come from the same point cloud. The network relates different positions within a single fractured piece to build a highly refined understanding of its own internal geometry and structural context.

\paragraph{Cross-Attention.} Once the internal geometry is understood, the network must find how two different pieces fit together. In cross-attention, the Queries come from the feature representations of Fragment A, while the Keys and Values come from Fragment B. This allows the points on one fractured surface to explicitly search the geometry of the other piece, facilitating the communication of local features across different fragments to find the exact point-to-point correspondences required for assembly.

\section{RANSAC}

The Random Sample Consensus (RANSAC) \cite{fischler1981random} algorithm is a robust mathematical estimation paradigm introduced by Fischler and Bolles. Rather than utilizing the entire dataset to compute an initial model and subsequently filtering out errors, RANSAC takes the opposite approach: it uses the smallest feasible subset of data to instantiate a model hypothesis, and then evaluates how much of the remaining data consistently aligns with that model.

\sloppy
Formally, given a dataset \(P\) and a model requiring \(n\) points to instantiate, RANSAC randomly selects a subset of n points to generate a hypothesis model, \(M1\). This model is then tested against all remaining points in P. Points that fall within a predefined error tolerance of M1 are classified as "inliers" and form the consensus set. If the size of this consensus set exceeds a specific threshold t, the model is accepted, and a refined model is typically recomputed using all identified inliers. If the consensus set is too small, the algorithm randomly samples a new subset of n points and repeats the process. After a predetermined number of iterations, the transformation that produced the largest consensus set is selected as the final solution.

The Jigsaw framework adopts this exact methodology to filter out false point-to-point correspondences predicted between fracture surfaces. Because neural networks inevitably predict some erroneous matches (outliers) that would severely degrade a global pose estimation, RANSAC is used to isolate the reliable geometric matches. In SE(3) space, the minimum number of points required to determine a rigid transformation is 3. Therefore, RANSAC randomly samples 3 corresponding point pairs from the two fragments.

To compute the transformation hypothesis from these 3 pairs, the algorithm utilizes Singular Value Decomposition (SVD) \cite{golub2013matrix}. First, the points from both fragments are centered at the origin by subtracting their respective centroids. These centered points are assembled into coordinate matrices \(A\) and \(B\). To find the optimal rotation matrix R that minimizes the least-squares error \(\lVert AR-B \rVert\), the algorithm computes the cross-covariance matrix \(H=A^TB\) and decomposes it via SVD such that \(H=U \sum V^T\). Because physical fragments must undergo a rigid transformation without any scaling or stretching, the singular values matrix \(\sum\) is discarded, and the pure rotation is extracted as \(R=UV^T\). The translation vector is then easily calculated by measuring the offset between the two original centroids after the rotation is applied.

Once this candidate transformation is calculated, it is applied to all other matched points to determine the consensus set. By repeating this hypothesis-and-evaluation loop, RANSAC successfully discards the hallucinated matches and outputs a highly reliable pairwise transformation based purely on the geometric inliers.
